% !TEX root = ../../../main.tex
% 来源：中学数学实验教材·第二册（几何）
% 章节：圆 / 与圆有关的角
% 原始文件：4.tex，行 2406-2451
% 类型：tikz
% ---
\begin{tikzpicture}
\begin{scope}
    \tkzDefPoints{0/0/O}
    \tkzDefPoint(70:1.5){A}
    \tkzDefPoint(-110:1.5){B}
    \tkzDefPoint(-30:1.5){C}
\tkzDrawCircle[thick](O,A)
\tkzDrawSegments[thick](A,B A,C)
\tkzDrawSegments[dashed](O,C)
\tkzAutoLabelPoints[center=O](A,B,C)
\tkzDrawPoints(O)
\tkzLabelPoints[left](O)
\node at (0,-2.3){(1)};
\tkzMarkAngles[mark=none, size=.4](B,O,C B,A,C A,C,O)

\end{scope}
\begin{scope}[xshift=4.5cm]
    \tkzDefPoints{0/0/O}
    \tkzDefPoint(70:1.5){A}
    \tkzDefPoint(-150:1.5){B}
    \tkzDefPoint(-110:1.5){D}
    \tkzDefPoint(-30:1.5){C}
\tkzDrawCircle[thick](O,A)
\tkzDrawSegments[thick](A,B A,C)
\tkzDrawSegments[dashed](A,D)
\tkzAutoLabelPoints[center=O](A,B,C,D)
\node at (0,-2.3){(2)};
\tkzLabelPoints[right](O)
\tkzDrawPoints(O)
\end{scope}
\begin{scope}[xshift=9cm]
    \tkzDefPoints{0/0/O}
    \tkzDefPoint(90:1.5){A}
    \tkzDefPoint(180:1.5){B}
    \tkzDefPoint(-90:1.5){D}
    \tkzDefPoint(-135:1.5){C}
\tkzDrawCircle[thick](O,A)
\tkzDrawSegments[thick](A,B A,C)
\tkzDrawSegments[dashed](A,D)
\tkzAutoLabelPoints[center=O](A,B,C,D)
\node at (0,-2.3){(3)};
\tkzDrawPoints(O)
\tkzLabelPoints[right](O)
\end{scope}

\end{tikzpicture}
